\documentclass[a4paper]{article}
\usepackage{graphicx}
\usepackage{amsmath,amssymb}
\usepackage{hyperref}
\usepackage{minted}
\usepackage{multirow}

\title{}
\date{}
\begin{document}
    \maketitle
    
    
    \section{Creational patterns}
        
        \paragraph{Factory method}
            \begin{minted}{python}
                from abc import ABC, abstractmethod
            
                class Document(ABC):
                    @abstractmethod
                    def create_page(self):
                            pass
            
                    class Resume(Document):
                        def create_page(self):
                            return "ResumePage"
            
                    class Report(Document):
                        def create_page(self):
                            return "ReportPage"
            \end{minted}
        
        \paragraph{Static factory method}
            
            \begin{minted}{python}
                class Car:
                    def init(self, brand):
                        self.brand = brand
                    @staticmethod
                    def toyota():
                        return Car("Toyota")
                my_car = Car.toyota()
            \end{minted}
        
        \paragraph{Abstarct factory}
        
        \paragraph{Simple factory}
            
            \begin{minted}{python}
                def shape_factory(kind):
                    if kind == "circle":
                        return Circle()
                    elif kind == "square":
                        return Square()
            \end{minted}
        
        \paragraph{Factory as a Function Object}
            
            \begin{minted}{python}
                class ConnectionFactory:
                    def call(self, host):
                        return Connection(host)
                factory = ConnectionFactory()
                conn = factory("localhost")
            \end{minted}
        
        \paragraph{Parameterized Factory}
            
            Factory methods that accept parameters to configure the returned object.
            
            Often combined with dependency injection.
        
        \paragraph{Virtual Constructor (a.k.a. Class Method Factory)}
            
            \begin{minted}{python}
                class Shape:
                    @classmethod
                    def from_sides(cls, sides):
                            if sides == 3:
                                return Triangle()
                            elif sides == 4:
                                return Square()
            \end{minted}
        
        \paragraph{Factory registry}
            
            The factory registry pattern is just a way to avoid a big if/elif chain by keeping a dictionary of “keys → constructor functions”. It’s basically a lookup table for creating objects.
            
            
            
            \begin{minted}{python}
                class AnimalFactory:
                    _registry: dict[str, type] = {}
                    @classmethod
                    def register(cls, name: str, animal_cls: type) -> None:
                        """Register a
                    class under a given name."""
                        cls._registry[name.lower()] = animal_cls
            
                    @classmethod
                    def create(cls, name: str, *args, **kwargs):
                        """Create an instance of the registered class."""
                        animal_cls = cls._registry.get(name.lower())
                        if animal_cls is None:
                            raise ValueError(f"Unknown animal type: {name}")
                        return animal_cls(*args, **kwargs)
            
                # Example classes
                class Dog:
                    def speak(self):
                        return "woof"
            
                class Cat:
                    def speak(self):
                        return "meow"
            
                # Register them
                AnimalFactory.register("dog", Dog)
                AnimalFactory.register("cat", Cat)
                # Usage
                a = AnimalFactory.create("dog")
                b = AnimalFactory.create("cat")
                print(a.speak()) # woof print(b.speak()) # meow
            \end{minted}
        
        \paragraph{Prototype}
            to build a copy from an existing object. So Cloning is the problem not using new. python has copy.deepcopy and shallowcopy for it
\end{document}